\subsubsection{Postmortem analiza iz core dump-a}\label{gdb-core-dump}

Kada program padne van debagera, nije uvek moguće ponoviti problem odmah. \emph{Core dump} čuva sliku procesa u trenutku pada: stekove niti, registre i delove memorije. To je osnovni format za postmortem analizu crash-eva.

Primer ima namernu grešku: funkcija \texttt{find\_user} može da vrati \texttt{NULL}, ali pozivalac to ne proverava.

\codefile{crash.c}{c}

Prevodimo program sa simbolima za debagovanje:

\begin{minted}{bash}
$ make
\end{minted}

Na sistemima koji dozvoljavaju core datoteke, jedan tok rada je:

\begin{minted}{bash}
$ ulimit -c unlimited
$ ./crash 42
Segmentation fault (core dumped)
$ gdb ./crash core
\end{minted}

Na distribucijama sa \texttt{systemd-coredump}, core datoteka često ne ostaje u radnom direktorijumu. Tada koristimo:

\begin{minted}{bash}
$ coredumpctl list crash
$ coredumpctl gdb crash
\end{minted}

U \texttt{gdb}-u prvo tražimo gde se program zaustavio:

\begin{minted}{text}
(gdb) bt
(gdb) bt full
(gdb) frame 0
(gdb) info locals
(gdb) p user
(gdb) frame 1
(gdb) p selected
(gdb) p wanted
\end{minted}

\texttt{bt} prikazuje stek poziva, a \texttt{bt full} dodaje lokalne promenljive za svaki okvir. U ovom primeru je neposredni pad u \texttt{print\_badge}, ali stvarni uzrok je prethodni okvir: \texttt{selected} je \texttt{NULL} jer traženi korisnik ne postoji.

Postmortem analiza dobro odgovara realnim incidentima: servis padne, sistem sačuva core, a analiza se radi kasnije na razvojnoj mašini sa istom binarnom verzijom i simbolima za debagovanje. Ključna disciplina je da se čuva tačna verzija izvršne datoteke, biblioteka i simbola koji odgovaraju core datoteci.
